% Generated draft: requires mathematical and source review.
\documentclass[11pt]{article}
\usepackage[margin=25mm]{geometry}
\usepackage{fontspec}
\setmainfont{Latin Modern Roman}
\usepackage{amsmath,amssymb,mathtools,unicode-math}
\setmathfont{Latin Modern Math}
\usepackage{enumitem,needspace,xurl,hyperref}
\setlength{\parindent}{0pt}
\setlength{\parskip}{5pt}
\setlength{\emergencystretch}{3em}
\newcommand{\N}{\mathbb{N}}
\newcommand{\Z}{\mathbb{Z}}
\newcommand{\Q}{\mathbb{Q}}
\newcommand{\R}{\mathbb{R}}
\newcommand{\C}{\mathbb{C}}
\newcommand{\F}{\mathbb{F}}
\newcommand{\A}{\mathbb{A}}
\newcommand{\PP}{\mathbb{P}}
\newcommand{\E}{\mathbb{E}}
\newcommand{\eps}{\varepsilon}
\newcommand{\dd}{\,\mathrm{d}}
\begin{document}
\section*{UnsolvedMath ID: 30003114}
\textbf{Source problem code:} catalogue entry
\subsection{Definitions and mathematical statement}
For each integer \(d\geq 0\), define
\[
\mathcal P_d=\left\{\sum_{j=0}^{d}a_jx^j\in\Z[x]:a_j\in\{-1,0,1\}\text{ for }0\leq j\leq d\right\}.
\]
Thus leading coefficients may vanish, and the zero polynomial belongs to \(\mathcal P_d\). The symbol \(\#S\) denotes the cardinality of a finite set \(S\).

A precise candidate reading of the supplied question is the following: does there exist a constant \(C>0\), independent of the rational number and of \(d\), such that, for every pair of coprime integers \(p,q\) with \(q>0\) and
\[
\frac{9}{10}\leq\frac pq\leq\frac{19}{20},
\]
there exists a constant \(c_{p,q}>0\), independent of \(d\), for which
\[
\forall d\in\Z_{\geq 1},\qquad
\#\left\{P\in\mathcal P_d:\left|P\!\left(\frac pq\right)\right|<c_{p,q}e^{-Cd}\right\}<e^{d/100}?
\]
Here both inequalities defining the small-value count and its upper bound are strict, and both endpoints of the rational interval are included.

\textit{Unresolved qualifications.} The displayed quantifier order interprets ``absolute'' as requiring one \(C\) for all rational numbers in the interval. The source places the existential phrase for \(C\) after the rational quantifier, so this uniformity requires confirmation. Positivity of \(c_{p,q}\), its independence of \(d\), the range \(d\geq1\), and the reduced-fraction convention are explicit interpretations rather than fully stated source hypotheses. If \(c_{p,q}>0\) and \(d=0\) is included, the zero polynomial makes the strict bound impossible, since its right-hand side is \(1\). The supplied background also describes this record as a conditional application of Question 4, but supplies neither that question nor the relevant conditional implication; no such implication can be stated precisely from this record alone.
\subsection{Short English statement}
The question asks whether one absolute exponential scale ensures that, at every rational point in \([9/10,19/20]\), fewer than \(e^{d/100}\) polynomials with coefficients in \(\{-1,0,1\}\) have extremely small values, with a positive prefactor allowed to depend on the rational point. The displayed formulation requires confirmation of the degree range, quantifier order, and missing connection to Question 4. This is a statement-only draft requiring human mathematical review.
\subsection{Sources}
\begin{itemize}
\item Supplied local UnsolvedMath registry record, permanent ID \(30003114\), exact source problem code \texttt{catalogue entry}, title ``Exponential Small-Value Bounds for Littlewood Polynomials,'' in the Oberwolfach Reports Open Problems collection. No individual proposer is identified; the supplied proposed year is 2016.
\item Supplied report citation: \emph{Diophantische Approximationen} (2017), Oberwolfach Reports, pp.~1099--1187, DOI \texttt{10.4171/owr/2016/21}: \url{https://doi.org/10.4171/owr/2016/21}. The record describes Oberwolfach Reports as documenting mathematical workshops and their open-problem discussions.
\item Supplied literature source, ``OWR report'': \url{https://ems.press/content/serial-article-files/46625?nt=1}.
\item Supplied literature source, ``Varju'': \url{https://arxiv.org/abs/1608.04210}. No fuller bibliographic description or precise supporting result is supplied.
\item Literature assessment dated exactly \texttt{2026-08-21}: the supplied curation labels the record open and describes it as a conditional application of the still-open Question 4, rather than a separately resolved theorem. This is supplied dated curation, not independently verified; the background expressly requests independent verification of time-sensitive status and cited claims.
\item The registry labels the statement \texttt{corrected\_verified} and says its mathematical content was preserved while notation, typography, or omitted definitions were repaired. These are supplied editorial labels, not independent mathematical verification of this draft.
\end{itemize}
\end{document}
